\subsubsection{The segment equation, derived}
\label{sec:front:segeq}

Equation~\eqref{eq:enth-ode} below is the heart of the march, and it is short
enough to look like a definition. It is not: it is the result of writing an
energy balance on two different control volumes and then eliminating the
quantity they share. This subsection does that in full, because the conductance
$K$ that comes out is \emph{not} the $\Ui P_{i}$ a reader would write down
unaided, and the difference reaches \SI{33}{\percent} at this design point.

\paragraph{The two control volumes.}
Fix a segment $j$, spanning $z\in[z_{j},z_{j+1}]$ with $\Delta z =
z_{j+1}-z_{j}$. Two control volumes share the tube wall over that interval:

\begin{description}[leftmargin=1.9em,itemsep=2pt]
\item[CV\,1 --- the PCM] belonging to segment $j$: cross-section
  $A_{\mathrm{avail}}$, length $\Delta z$, stationary, closed.
\item[CV\,2 --- the fluid] inside the tube over the same interval:
  cross-section $A_{\mathrm{flow}}=\pi r_{i}^{2}$, open, with mass flow $\md$
  entering at $T_{j}$ and leaving at $T_{j+1}$.
\end{description}

\noindent
Write $q'_{j}$ for the heat per unit tube length crossing from CV\,2 into
CV\,1. It appears in both balances with opposite sign, and eliminating it is
the whole exercise.

\medskip
\noindent\textbf{Step 1 --- the PCM balance.}
CV\,1 is closed and stationary, so the first law reduces to a statement about
its enthalpy alone. With $E'$ the enthalpy per unit tube length
(\S\ref{sec:front:enthalpy}), the enthalpy of the control volume is
$E'\Delta z$ and
\begin{equation}
  \frac{\mathrm{d}}{\mathrm{d}t}\bigl(E'_{j}\,\Delta z\bigr)
  \;=\;
  \underbrace{q'_{j}\,\Delta z}_{\text{through the tube wall}}
  \;\underbrace{-\;\bigl[\dot Q_{\mathrm{ax}}\bigr]_{z_{j}}^{z_{j+1}}}
              _{\text{axial conduction, H2}}
  \;+\;\underbrace{\dot Q_{\mathrm{form}}}_{\text{formation, H6}}
  \;+\;\underbrace{\dot Q_{\mathrm{azi}}}_{\text{neighbours, H9}} ,
\end{equation}
and the last three terms are dropped by hypothesis. Two of those deletions are
worth a number rather than a citation:

\begin{itemize}[leftmargin=1.5em,itemsep=2pt]
\item \textbf{Axial conduction (H2)} is negligible by an enormous margin. The
  axial diffusion time over one segment is $\Delta z^{2}/\alpha_{l} =
  \SI{5.4e9}{\second}$ against a \SI{3.6e4}{\second} charging window --- a ratio
  of $1.5\times10^{5}$. With $\Delta z = \SI{30.5}{\meter}$ this is not a close
  call, and it would remain safe at a hundred times the segment count.
\item \textbf{Azimuthal conduction (H9)} is \emph{not} negligible and is
  assumed away pending the calculation described in \S\ref{sec:defects:h9}. At
  the wellhead the neighbouring cell is \SI{48.9}{\kelvin} away and a slab
  estimate gives $\sim\SI{22}{\watt\per\meter}$ against a useful duty of
  $\sim\SI{83}{\watt\per\meter}$. This is the least defensible term in the
  model.
\end{itemize}

\noindent
Dividing by $\Delta z$ leaves the PCM balance in the form used by the code:
\begin{equation}
  \boxed{\;\frac{\mathrm{d}E'_{j}}{\mathrm{d}t} \;=\; q'_{j}\;}
  \label{eq:pcm-balance}
\end{equation}

\medskip
\noindent\textbf{Step 2 --- the fluid balance.}
CV\,2 is open. For incompressible single-phase flow (H7) with no viscous
dissipation, the energy equation per unit tube length is
\begin{equation}
  \underbrace{\rho A_{\mathrm{flow}} c_{p}\,\frac{\partial T}{\partial t}}
             _{\text{fluid thermal inertia}}
  \;+\;
  \md\,c_{p}\,\frac{\partial T}{\partial z}
  \;=\;
  -\,q'(z) .
  \label{eq:fluid-pde}
\end{equation}

The model drops the first term, and this deserves justification rather than
assertion. Two comparisons, pointing the same way:

\begin{center}
\begin{tabular}{@{}l r l@{}}
\toprule
fluid transit time over $L_{\mathrm{tube}}$ & \SI{2264}{\second} & ($u=\SI{1.35}{\meter\per\second}$) \\
charging window $t_{\mathrm{ch}}$           & \SI{36000}{\second} & \\
\midrule
ratio                                        & \num{0.063} & \\
$\rho A_{\mathrm{flow}}c_{p} \,/\, C_{l}$    & \num{0.257} & fluid vs.\ PCM sensible capacity \\
$\rho A_{\mathrm{flow}}c_{p}\Delta T \,/\, E'_{\mathrm{lat}}$
   & \num{0.0122} & fluid swing vs.\ PCM latent, $\Delta T=\SI{10}{\kelvin}$ \\
\bottomrule
\end{tabular}
\end{center}

\noindent
The fluid re-establishes its profile in one transit, $\SI{2264}{\second}$,
while the PCM state evolves over \SI{10}{\hour}: a $16\times$ separation, so
quasi-steady is reasonable but \emph{not} asymptotic. The honest reading is
that the first \SI{38}{\minute} of each half-cycle is a startup transient the
model does not resolve, worth about \SI{6}{\percent} of the window. What makes
this tolerable is the third row: the energy parked in the fluid is only
\SI{1.2}{\percent} of the latent inventory it is charging, so the transient
misplaces very little energy even though it lasts a noticeable fraction of the
run.

Setting $\partial T/\partial t = 0$ and using the fact that CV\,1 presents a
\emph{single} temperature $T_{\mathrm{pcm}}$ to the whole segment
(H5, the lumped state of \S\ref{sec:front:enthalpy}), the driving difference is
$T(z)-T_{\mathrm{pcm}}$ and the local flux closes on the network of
\S\ref{sec:network}:
\begin{equation}
  q'(z) \;=\; \Ui\,\bigl(2\pi r_{i}\bigr)\,\bigl[T(z)-T_{\mathrm{pcm}}\bigr].
  \label{eq:local-flux}
\end{equation}
Note that $2\pi r_{i}$ appears, not $P_{T}$: $\Ui$ was referred to the
\emph{inner} area in Eq.~\eqref{eq:Ui}, and the finned perimeter is already
inside it through $R'_{\mathrm{outer}}$. Using $P_{T}$ here would count the fins
twice --- which is a compact statement of the v0.1 defect.

\medskip
\noindent\textbf{Step 3 --- integrate across the segment.}
Substituting Eq.~\eqref{eq:local-flux} into the steady form of
Eq.~\eqref{eq:fluid-pde} gives a linear ODE in $z$ with $T_{\mathrm{pcm}}$ and
$\Ui$ held constant over the segment:
\begin{equation}
  \md c_{p}\frac{\mathrm{d}T}{\mathrm{d}z}
  = -2\pi r_{i}\Ui\bigl(T-T_{\mathrm{pcm}}\bigr)
  \quad\Longrightarrow\quad
  \frac{\mathrm{d}\bigl(T-T_{\mathrm{pcm}}\bigr)}{T-T_{\mathrm{pcm}}}
  = -\frac{2\pi r_{i}\Ui}{\md c_{p}}\,\mathrm{d}z .
\end{equation}
Integrating from $z_{j}$ to $z_{j+1}$,
\begin{equation}
  T_{j+1} \;=\; T_{\mathrm{pcm}}
    + \bigl(T_{j}-T_{\mathrm{pcm}}\bigr)e^{-\mathrm{NTU}_{j}},
  \qquad
  \mathrm{NTU}_{j} \;=\; \frac{2\pi r_{i}\,\Ui\,\Delta z}{\md\,c_{p}},
  \label{eq:ntu-derived}
\end{equation}
which is Eq.~\eqref{eq:ntu}. It is the single-stream heat-exchanger result: one
stream with finite capacity rate $\md c_{p}$, the other an isothermal reservoir
of infinite capacity rate, effectiveness $\varepsilon = 1-e^{-\mathrm{NTU}}$.
The reservoir idealisation is exactly what the lumped PCM state licenses.

\medskip
\noindent\textbf{Step 4 --- the closure.}
The heat the fluid actually gives up over the segment is the enthalpy it loses,
\begin{equation}
  q'_{j}\,\Delta z \;=\; \md c_{p}\bigl(T_{j}-T_{j+1}\bigr)
  \;\overset{\eqref{eq:ntu-derived}}{=}\;
  \md c_{p}\bigl(T_{j}-T_{\mathrm{pcm}}\bigr)\bigl(1-e^{-\mathrm{NTU}_{j}}\bigr),
\end{equation}
and dividing by $\Delta z$ gives the conductance that closes
Eq.~\eqref{eq:pcm-balance}:
\begin{equation}
  \boxed{\;
  q'_{j} = K_{j}\bigl(T_{j}-T_{\mathrm{pcm}}\bigr),
  \qquad
  K_{j} = \frac{\md\,c_{p}\bigl(1-e^{-\mathrm{NTU}_{j}}\bigr)}{\Delta z} .
  \;}
  \label{eq:K-derived}
\end{equation}

Combining Eqs.~\eqref{eq:pcm-balance} and \eqref{eq:K-derived} with the
enthalpy curve $T_{\mathrm{pcm}}(E')$ of Eq.~\eqref{eq:enth-curve} closes the
system, and the result is Eq.~\eqref{eq:enth-ode}.

\paragraph{Why $K$ is not $\Ui P_{i}$.}
This is the step most easily got wrong, so it is worth isolating. The naive
closure takes Eq.~\eqref{eq:local-flux} with $T(z)\approx T_{j}$, giving
$q'_{j}\approx 2\pi r_{i}\Ui (T_{j}-T_{\mathrm{pcm}})$. Comparing:
\begin{equation}
  \frac{K_{j}}{2\pi r_{i}\Ui}
  \;=\;
  \frac{1-e^{-\mathrm{NTU}_{j}}}{\mathrm{NTU}_{j}}
  \;\le\; 1 .
  \label{eq:K-ratio}
\end{equation}
The two limits say what the difference means:

\begin{description}[leftmargin=1.9em,itemsep=3pt]
\item[$\mathrm{NTU}\to0$:] the ratio $\to1$. The fluid barely cools across the
  segment, $T(z)\approx T_{j}$ throughout, and the naive closure is right.
\item[$\mathrm{NTU}\to\infty$:] $K\to \md c_{p}/\Delta z$, a finite limit set
  by the \emph{fluid}, not the wall. No matter how good the borehole is, a
  segment cannot absorb more than the stream can carry to it. The naive closure
  has no such limit and would let $T_{j+1}$ fall below $T_{\mathrm{pcm}}$ ---
  heat flowing from cold to hot, at a rate that grows without bound as $\Ui$
  improves.
\end{description}

\noindent
At this design point $\mathrm{NTU}$ runs from \num{0.083} to \num{0.602} over
the charge, so Eq.~\eqref{eq:K-ratio} lies between \num{0.96} and \num{0.75}:
\textbf{using $\Ui P_{i}$ would overstate the conductance by \SI{4}{\percent}
to \SI{33}{\percent}}, worst where the melt layer is thinnest and the borehole
best. The error is largest exactly where the model is most active, which is why
it cannot be waved away as second order.

\paragraph{Two properties inherited from this derivation.}

\begin{description}[leftmargin=1.9em,itemsep=3pt]
\item[Closure is structural.] The code does not evaluate
  Eq.~\eqref{eq:K-derived} and Eq.~\eqref{eq:pcm-balance} independently. It
  integrates Eq.~\eqref{eq:pcm-balance} over the step and \emph{reports}
  $q'_{j}=(E'^{\,n+1}-E'^{\,n})/\Delta t$, then drops the fluid temperature by
  exactly that amount, $T_{j+1}=T_{j}-q'_{j}\Delta z/(\md c_{p})$. Whatever the
  PCM gained, the fluid lost, to machine precision, with no possibility of
  drift.
\item[$K$ is frozen over a time step, $T_{\mathrm{pcm}}$ is not.] $K_{j}$ is
  evaluated once per step from the melt thickness at the start of it, whereas
  $T_{\mathrm{pcm}}$ is advanced exactly inside the step by
  Eq.~\eqref{eq:enth-exact}. The consequence is that $T_{j+1}$ equals the
  $\varepsilon$--NTU result of Eq.~\eqref{eq:ntu-derived} only while the segment
  is on the latent plateau, where $T_{\mathrm{pcm}}=\Tm$ is genuinely constant.
  In the sensible branches the two differ at second order in $\Delta t$. This
  is a deliberate choice: it keeps closure exact, which matters more than
  matching an $\varepsilon$--NTU relation that is itself only quasi-steady.
\end{description}
